\section*{Ukázka:}

V~hospodě U~kalicha seděl jen jeden host. Byl to civilní strážník Bretschneider, stojící ve službách státní policie. Hostinský Palivec myl tácky a~Bretschneider se marně snažil navázat s~ním vážný rozhovor.

Palivec byl známý sprosťák, každé jeho druhé slovo byla zadnice nebo hovno. Přitom byl ale sečtělý a~upozorňoval každého, aby si přečetl, co napsal o~posledním předmětě Viktor Hugo, když líčil poslední odpověď staré gardy Napoleonovy Angličanům v~bitvě u~Waterloo.

„To máme pěkné léto,” navazoval Bretschneider svůj vážný rozhovor.

„Stojí to všechno za hovno,” odpověděl Palivec, ukládaje tácky do skleníku.

„Ty nám to pěkně v~tom Sarajevu vyvedli,” se slabou nadějí ozval se Bretschneider.

„V~jakým Sarajevu?” otázal se Palivec, „v tej nuselskej vinárně? Tam se perou každej den, to vědí, Nusle.”

„V~bosenském Sarajevu, pane hostinský. Zastřelili tam pana arcivévodu Ferdinanda. Co tomu říkáte?”

„Já se do takových věcí nepletu, s~tím ať mně každej políbí prdel,” odpověděl slušně pan Palivec, zapaluje si dýmku, „dneska se do toho míchat, to by mohlo každému člověkovi zlomit vaz. Já jsem živnostník, když někdo přijde a dá si pivo, tak mu ho natočím. Ale nějaký Sarajevo, politika nebo nebožtík arcivévoda, to pro nás nic není, z~toho nic nekouká než Pankrác.”

Bretschneider umlkl a~díval se zklamaně po pusté hospodě. „Tady kdysi visel obraz císaře pána,” ozval se opět po chvíli, „právě tam, kde teď visí zrcadlo.”

„Jó, to mají pravdu,” odpověděl pan Palivec, „visel tam a~sraly na něj mouchy, tak jsem ho dal na půdu. To víte, ještě by si někdo mohl dovolit nějakou poznámku a~mohly by být z~toho nepříjemnosti. Copak to potřebuju?”

„V~tom Sarajevu to muselo být asi ošklivý, pane hostinský.”

Na tuto záludně přímou otázku odpověděl pan Palivec neobyčejně opatrně: „V tuhle dobu bývá v~Bosně a~Hercegovině strašný horko. Když jsem tam sloužil, tak museli dávat našemu obrlajtnantovi led na hlavu.”

„U kterého pluku jste sloužil, pane hostinský?”

„Na takovou maličkost se nepamatuju, já jsem se nikdy o~takovou hovadinu nezajímal a~nikdy jsem nebyl na to zvědavej,” odpověděl pan Palivec, „přílišná zvědavost škodí.”

Civilní strážník Bretschneider definitivně umlkl a~jeho zachmuřený výraz se zlepšil teprve příchodem Švejka, který, vstoupiv do hospody, poručil si černé pivo s~touto poznámkou: „Ve Vídni dneska taky mají smutek.”

\cite{osudy_dobreho_vojaka_svejka}

\newpage

